\section{Completion of the proof}

\begin{proof}[Proof of Theorem~\ref{thm:main}]
For \(\sqrt3\le c\le2\), the theorem is
\cite[Theorem~2.8, with Definitions~2.6--2.7]{DKLM2026}, in precisely the
unit-height and Green-kernel conventions used here. The source's
regularity estimates also retain the lattice cutoff when distinct
continuum points round to the same face. Thus this interval includes
the endpoint \(c=2\).

Fix \(0<c<\sqrt3\). Sections~\ref{sec:finite}
and~\ref{sec:thermodynamic} construct the limits of every finite word
of rows and marks. Proposition~\ref{prop:physical-plane} identifies the
physical words with local arrow probabilities of the prescribed plane
law. For \(c<1\) the construction uses the exactly balanced sector
throughout and proves existence in the required iterated order. Where
an auxiliary near-balanced limit is used for \(c\ge1\), the equilibrium
identification gives the already existing balanced law. The same
proposition identifies the two axial descriptions of its local
observables.

The angular representation of Section~\ref{sec:angular} has orbit mass
\(\kappa<\infty\) and height variogram
\(4\kappa\log|x|+O(1)\). Proposition~\ref{prop:gaussian-currents}
establishes the holomorphic and antiholomorphic Wick rules, with
logarithmic covariance \(-2\kappa\log|x-y|\).
Propositions~\ref{prop:smooth-tests}, \ref{prop:negative-regularity},
\ref{prop:offdiagonal}, and~\ref{prop:finite-energy} then give all the
claimed convergences with squared GFF multiplier \(4\pi\kappa\).
Their estimates apply to arbitrary real mesh sizes and identify every
subsequential limit; hence the limits hold along the full
\(\delta\downarrow0\).

If \(0<c<1\), Proposition~\ref{prop:normalization} gives
\[
 4\pi\kappa=\frac2{\pi-\lambda}=\frac2{\arccos\Delta}.
\]
For \(1\le c<\sqrt3\), we have now proved a subsequential GFF limit in
the required senses for the balanced law, so the hypothesis of
\cite[Theorem~4.3]{DKLM2026} is satisfied. That theorem identifies its
squared multiplier as \(-1/f''(0)\), where \(f\) is the source's
free-energy function. By \cite[Theorem~4.4]{DKLM2026},
\(f''(0)=-\tfrac12\arccos\Delta\), giving the same multiplier.
This use of the conditional theorem comes after, and does not supply,
the convergence argument.

Finally, ordinary finite signed-energy tests have the joint covariance
in the statement by Proposition~\ref{prop:finite-energy}. All pinning
choices differ by an additive constant, the negative-regularity assertions
use the common chosen pin, and increments cancel it. The uniformity in
\eqref{eq:increment-main} and the arbitrary-open-set conclusions have
already been proved in their respective propositions. This completes
the theorem.
\end{proof}
